\begin{lemma}
\label{editorial-source-lemma-conclude-jacobson-Noetherian}
With notation as above, assume that $R$ is a Noetherian Jacobson ring.
Further assume $R \to S$ is of finite type.
There is a commutative diagram
$$
\xymatrix{
\text{Constr}(Y) \ar[r]^{E \mapsto E_0} \ar[d]^{E \mapsto f(E)} &
\text{Constr}(Y_0) \ar[d]^{E \mapsto f(E)} \\
\text{Constr}(X) \ar[r]^{E \mapsto E_0} &
\text{Constr}(X_0)
}
$$
where the horizontal arrows are the bijections from
Topology, Lemma \ref{topology-lemma-jacobson-equivalent-constructible}.
The same diagram and conclusions hold more generally when $R$ is
Jacobson and $R\to S$ is of finite presentation, without requiring
$R$ to be Noetherian.
\end{lemma}

\begin{proof}
Under the original hypotheses, the finite type map $R\to S$ is of
finite presentation by
Lemma \ref{editorial-source-lemma-Noetherian-finite-type-is-finite-presentation}.
For the more general assertion, finite presentation is assumed directly.
In both cases it implies finite type, so $S$ is Jacobson by
Proposition \ref{editorial-source-proposition-Jacobson-permanence}. Therefore $X$ and
$Y$ are Jacobson spaces and the two horizontal maps are bijections
by Topology, Lemma \ref{topology-lemma-jacobson-equivalent-constructible},
which requires no Noetherian hypothesis.
Chevalley's theorem (Theorem \ref{editorial-source-theorem-chevalley}) states that
the image of a constructible subset of the source $Y$ is constructible
in the target $X$. For each such $E\subset Y$,
Lemma \ref{editorial-source-lemma-image-finite-type-map-Jacobson-rings} gives
$f(E_0)=f(E)_0$, proving commutativity.
The right vertical arrow is well defined as a direct-image map:
given a constructible subset $C\subset Y_0$, the top horizontal
bijection supplies a unique constructible $E\subset Y$ with $E_0=C$,
and then $f(C)=f(E)_0$ is constructible in $X_0$ by the bottom
horizontal bijection. This proves every arrow and the square in
the stated greater generality.
\end{proof}